% ===========================================================================
% PART A -- WHAT NORMATIVE INFORMATION IS. The concept, examined exhaustively.
% This is the intellectual core and it was previously scattered across the paper.
% ===========================================================================
\section{What normative information is}
\label{sec:concept}

The phrase appears in every paper on value alignment and is almost never defined.
This section defines it, by killing five candidate definitions and keeping the
sixth. Each candidate is a position someone actually holds; each is refuted by a
specific construction, not by preference.

\subsection{The naive picture, and precisely where it breaks}

The naive picture is a pipe. A person has values; they write them down; the
writing is aggregated; the aggregate is compiled; a model is trained. At each
stage some fraction of the original content is lost, and one could in principle
report that fraction.

\step{Where it breaks.} ``Fraction of \emph{what}'' has no answer. Values are not
a substance with a quantity. Before any measurement is possible one must say what
two normative objects being ``the same'' means, and that is a modelling decision
with at least six defensible answers.

\subsection{Candidate 1 --- normative information is the text}

\emph{Claim.} What a person wrote is what they contributed; loss is text that
failed to survive.

\begin{proposition}[Refuted by paraphrase]
Under this definition, rewriting a criterion in different words with identical
meaning is total loss, and copying a criterion verbatim into a context where it
cannot apply is zero loss. Both readings are wrong in the direction that matters
for a standard.
\end{proposition}

\step{The empirical form of the refutation here.} The compiled standard rewrites
every criterion. Under Candidate 1 the compilation destroys everything, which is
plainly false: the compiled standard still selects the same response as the full
rubric on $70.6\%$ of prompts (§\ref{sec:compile}). A definition under which a
$70.6\%$-agreeing object has lost $100\%$ of the content is not measuring content.

\subsection{Candidate 2 --- normative information is the numbers}

\emph{Claim.} The weights are the content; loss is weights that changed.

\begin{proposition}[Refuted by the units problem]
The weights live on a scale with no established zero for comparison across
people, no established unit, and --- as §\ref{sec:cardinal} measures --- an
interval reading that changes $11.7\%$ of decisions when removed. A definition
that treats an uncalibrated private scale as the content inherits every one of
those problems as though they were facts about values.
\end{proposition}

\step{And a sharper refutation.} Multiply everyone's weights by two. Every
decision is unchanged, every ranking is unchanged, every aggregate is unchanged.
Under Candidate 2 all the information changed. So Candidate 2 measures a
representation, not a content.

\subsection{Candidate 3 --- normative information is the induced ranking}

\emph{Claim.} What matters is the order over the candidates; loss is order that
changed.

This is much better, and it is what most evaluation work implicitly uses. It is
still wrong, in two ways that the data exhibit.

\begin{proposition}[Refuted by candidate-set dependence]
\label{prop:cand3}
An ordering is defined only relative to the set being ordered. §\ref{sec:agg}
measures that $13.4\%$ of pairwise relations \emph{reverse} when an irrelevant
third candidate is removed. So ``the induced ranking'' is not a property of the
person; it is a property of the person and the four responses shown.
\end{proposition}

\begin{proposition}[Refuted by type]
\label{prop:cand3b}
Two normative objects can induce the same ranking on every candidate set yet
differ in a way that matters: one says ``$D$ is worse'', the other says ``$D$ is
forbidden''. On any set where $D$ does not win they are indistinguishable by
ranking. \Cref{thm:veto} shows this difference is not recoverable by any finite
weight, and §\ref{sec:agg} measures that it bites: $26.1\%$ of the time the
aggregate selects a response somebody forbade.
\end{proposition}

\subsection{Candidate 4 --- normative information is what a model trained on it does}

\emph{Claim.} The only thing that matters is downstream behaviour.

\begin{proposition}[Not refuted, but not available]
This is a coherent definition. It requires $Y$: the behaviour of a system trained
under the standard. $Y$ is not in the release and no arithmetic recovers it
(§\ref{sec:notid}). Any number quoted end-to-end under this definition is
fabricated.
\end{proposition}

\begin{remark}[Why it is also unstable as a definition]
Even with $Y$, the answer would depend on the training method, the base model,
the optimiser, and the data mixture. The measurement would then be about the
training stack, and comparing two standards would require holding all of it
fixed --- which is a legitimate experiment and a different one from measuring the
standards.
\end{remark}

\subsection{Candidate 5 --- normative information is mutual information}

\emph{Claim.} Use Shannon: $I(\Nn; G)$, the mutual information between the
individual norms and the aggregate.

\begin{proposition}[Refuted by decision-irrelevance]
Mutual information counts \emph{any} statistical dependence, including
dependence that no decision uses. A compilation that faithfully preserves a
criterion's exact wording while destroying its polarity scores high mutual
information and is catastrophic; one that discards wording while preserving every
decision scores low and is harmless. Shannon information is the wrong currency
because it has no notion of what the information is \emph{for}.
\end{proposition}

\begin{remark}[The rate--distortion repair, and its remaining hole]
The standard fix is rate--distortion: minimise bits subject to a distortion
bound. §\ref{sec:compile} runs exactly this and gets a real answer ($1$ bit
reproduces $89.8\%$ of decisions). But the answer is entirely determined by the
choice of distortion measure, which is the same freedom Candidate 6 makes
explicit. Rate--distortion is not an alternative to declaring the purpose; it is
a way of computing once the purpose is declared.
\end{remark}

\subsection{Candidate 6 --- normative information is a query-relative equivalence class}

\begin{definition}[The definition this paper uses]
\label{def:normativeinfo}
Fix a family $Q$ of \emph{queries}: the questions the downstream system must be
able to answer. Then
\[
\Nn \sim_Q \Nn' \iff q(\Nn,x)=q(\Nn',x) \ \text{ for all } q\in Q,\ \text{all } x,
\]
and the \emph{normative information} in $\Nn$ is the equivalence class $[\Nn]_Q$.
\end{definition}

\step{Why this is not a dodge.} It makes the freedom that Candidates 1--5 each
exercised silently into an explicit, publishable input. Candidate 1 is
\cref{def:normativeinfo} with $Q$ = all text functions; Candidate 3 is $Q$ = the
ranking queries on the observed set; Candidate 5 is $Q$ = all measurable
functions. They are not rival definitions --- they are different choices of $Q$
inside one definition, and \cref{prop:Q} proves the choice determines the answer.

\mesotakeaway{\textbf{Declaring $Q$ is not a preliminary to the measurement.
Declaring $Q$ \emph{is} most of the result.} Any framework reporting a single
percentage of normative information preserved has chosen a $Q$ and not told you.}

\subsection{Where the definition lives formally: Blackwell}

\Cref{def:normativeinfo} is a statement about indistinguishability. Decision
theory has the matching notion.

\begin{itemize}
\item Take $Q$ to be a family of \emph{decision problems} rather than arbitrary
      questions --- which loses no generality, because any query can be posed as
      a decision.
\item Then $\Nn \sim_Q \Nn'$ exactly when neither is \emph{more informative} than
      the other in Blackwell's sense on that family.
\item And the natural quantity is not a percentage but a \emph{deficiency}:
      $\delta_Q(\Nn,\Nn') = \sup_{d\in Q}[R_{\Nn'}(d)-R_{\Nn}(d)]$, the worst
      extra loss over the declared family.
\end{itemize}

\begin{remark}[What this buys, concretely]
Three things a percentage cannot give. (i) A \emph{unit}: the deficiency is in
the units of the decision loss, so it can be compared to what the decision is
worth. (ii) A \emph{worst case}: a supremum over the family, not an average over
whatever queries happened to be tried. (iii) An \emph{ordering}: Blackwell's
relation is a partial order, so two standards can be genuinely incomparable, and
saying so is more honest than forcing a ranking.
\end{remark}

\subsection{The three source objects, and why merging them makes the question unanswerable}

\begin{definition}[Layers]
\label{def:layers}
\begin{itemize}
\item $Z = (\Nn_1,\dots,\Nn_n)$: the individual inputs. Who said what, who
      dissented, who vetoed.
\item $G = \Agg(Z)$: the collective norm, under a rule $\Agg$.
\item $S$: the authorised behaviour set, after any legitimate expert enrichment.
\end{itemize}
\end{definition}

\begin{proposition}[Without the layers, ``defect'' is undefined]
Aggregation necessarily destroys individual information (\cref{thm:garble}).
Therefore an observed loss between $Z$ and $G$ is a \emph{defect} only if it was
not licensed by $\Agg$. If $\Agg$ is not declared, no observed loss can be
classified, and ``the pipeline lost information'' cannot be distinguished from
``the pipeline did what it said it would''.
\end{proposition}

\step{The state of the object.} CoVal declares no $\Agg$. So at this site the
distinction between an authorised social choice and a machine defect is
\textbf{not recoverable}. This is not a hard measurement problem; it is a missing
line of documentation, and it is the single cheapest thing a future release could
fix.

\subsection{The four classes of distinction, which is what a conservation law looks like}

The goal of a theory of normative information is not a score. It is a statement
of \emph{what may be compressed and what may not} --- the normative analogue of a
conservation law, where the content is the partition rather than the total.

\begin{table}[H]\centering\small
\begin{tabularx}{\textwidth}{P{3.4cm}P{5.2cm}Y}
\toprule
class & distinctions & what is required of the pipeline\\
\midrule
\textbf{must be causally preserved} & polarity; non-compensability (veto); scope;
exception conditions & intervening on these must move behaviour, in the predicted
direction\\
\textbf{recoverable is enough} & wording; provenance where not load-bearing for
authority; redundant restatements & must be decodable; need not steer\\
\textbf{licensed to drop} & duplicates; criteria entailed by others; phrasing
variance & only if the licence is declared at the layer where it happens\\
\textbf{type errors when lost} & veto $\to$ scalar; uncertainty $\to$ certainty;
right $\to$ preference; expert constraint $\to$ public preference & not
degradation but a category change; no continuous metric detects it\\
\bottomrule
\end{tabularx}
\caption{The partition this research programme is aimed at filling empirically.
Rows two and three are currently assertions. Row one is partly measured
(§\ref{sec:est}). Row four is where this paper's theorem sits.}
\label{tab:conservation}
\end{table}

\subsection{The deepest distinction: type versus magnitude}

Rows one and four of \cref{tab:conservation} are the ones that make the framework
falsifiable, and the fourth is the deepest, because it is the only one that
\emph{no similarity metric can detect}.

\step{Why no metric detects it.} A type change can leave the object arbitrarily
close to the source in every continuous measure while answering a different
question. A veto rendered as $-10$ and a strong dispreference rendered as $-10$
are \emph{identical objects}. There is nothing left to measure. The difference
exists only in the behaviour under candidate sets not yet seen --- which is
exactly \cref{thm:veto}, and exactly why that theorem is the load-bearing result
of this paper.

\mesotakeaway{Normative information is not a quantity that flows down a pipe. It
is a \textbf{partition of distinctions into what must be causally preserved, what
must merely be recoverable, what may be dropped under a declared licence, and
what cannot be represented at all in the chosen format}. The last class is
invisible to every similarity measure, and it is where a scalar scale fails.}
